%--- E4: end of Section 3, before the keypoint ------------------------------
\begin{explained}{the ODE prototype, step by step}
This section is the one-dimensional version of the whole theory. It starts with an ODE because in one variable the algebra is completely explicit: roots of a polynomial determine the solutions, and the roots lead naturally to the minimal annihilator, the companion matrix and the spectral projectors.

\medskip\textbf{Step 1: the ODE.} $p(D)u=0$ with $D=d/dx$ and $p$ of degree $m$, factored over $\C$ as $p(\lambda)=\prod_{j=1}^r(\lambda-\lambda_j)^{m_j}$, $\sum_jm_j=m$: the $\lambda_j$ are the distinct roots, $m_j$ their multiplicities, $m$ the total degree. For $p(\lambda)=(\lambda-1)^2(\lambda+2)^3$ the equation $(D-1)^2(D+2)^3u=0$ has order five.

\medskip\textbf{Step 2: why solutions are $q(x)e^{\lambda x}$.} First, $(D-\lambda)e^{\lambda x}=0$. If the root has multiplicity greater than one, further solutions appear: $(D-\lambda)(xe^{\lambda x})=e^{\lambda x}$, hence $(D-\lambda)^2(xe^{\lambda x})=0$; in general $(D-\lambda)^m(x^ke^{\lambda x})=0$ for $0\le k<m$. The reason is the identity
\[
(D-\lambda)\big(q(x)e^{\lambda x}\big)=q'(x)e^{\lambda x}:
\]
each application of $D-\lambda$ differentiates the polynomial prefactor once, and after $m$ applications a polynomial of degree $<m$ has vanished. A root of multiplicity $m_j$ therefore contributes the $m_j$ solutions $e^{\lambda_jx},xe^{\lambda_jx},\dots,x^{m_j-1}e^{\lambda_jx}$, and the complete solution space is
\[
\ker p(D)=\Big\{\sum_{j=1}^rq_j(x)e^{\lambda_jx}:\deg q_j<m_j\Big\}.
\]
\emph{Worked example.} For $(\lambda-1)^2(\lambda+2)^3$: root $1$ (multiplicity 2) gives $e^x,xe^x$; root $-2$ (multiplicity 3) gives $e^{-2x},xe^{-2x},x^2e^{-2x}$. Five independent solutions, so the solution space has dimension five $=m$: the number of natural basis functions is finite and equals the degree of $p$.

\medskip\textbf{Step 3: the minimal annihilator.} For a particular function $u$, $\Ann(u)=\{Q\in\C[D]:Q(D)u=0\}$. In one variable the annihilator of a nonzero $u$ is generated by a single polynomial of smallest degree; normalized to be monic it is the \emph{minimal annihilating polynomial} of $u$. Example: $u=xe^x$ has $(D-1)u=e^x$ and $(D-1)^2u=0$, so its minimal annihilator is $p_u(D)=(D-1)^2$; the multiplicity two tells that the solution has a polynomial prefactor of degree one. The polynomial $p$ defining the ODE need not be minimal for a particular solution: $u=e^x$ solves $(D-1)^2u=0$, but its minimal annihilator is only $D-1$.

\medskip\textbf{Step 4: the companion matrix.} Write $p(\lambda)=\lambda^m+a_{m-1}\lambda^{m-1}+\dots+a_0$. Then $u^{(m)}=-a_{m-1}u^{(m-1)}-\dots-a_0u$. With the state vector $Y=(u,u',\dots,u^{(m-1)})^\top$ the ODE becomes $Y'=CY$, with $C$ the companion matrix displayed in the text. Its characteristic polynomial is $p(\lambda)$: the roots of the polynomial become the eigenvalues of a matrix, and the multiplicities determine the possible Jordan blocks.

\medskip\textbf{Step 5: spectral projectors.} Suppose first the roots are distinct. The Lagrange polynomials $\ell_j(z)=\prod_{k\ne j}\frac{z-\lambda_k}{\lambda_j-\lambda_k}$ satisfy $\ell_j(\lambda_i)=1$ if $i=j$ and $0$ if $i\ne j$. Substituting the matrix gives the projectors $\Pi_j=\ell_j(C)$; each projects onto the eigenspace of $\lambda_j$, and $\Pi_i\Pi_j=0$ for $i\ne j$, $\sum_j\Pi_j=\Id$. Consequently $e^{xC}=\sum_je^{\lambda_jx}\Pi_j$ when $C$ is diagonalizable with distinct eigenvalues. When roots are repeated, eigenspace projectors alone do not capture the polynomial factors: one needs the generalized eigenspaces and the nilpotent parts of the Jordan blocks, which produce $xe^{\lambda x}$, $x^2e^{\lambda x}$, \dots

\medskip\textbf{Step 6: why this is the prototype.}
\begin{center}\small
\begin{tabular}{@{}L{5.6cm}L{8.6cm}@{}}
\toprule
One-dimensional ODE & Later multivariable theory\\\midrule
polynomial $p(\lambda)$ & polynomial symbol $P(\tau,\xi)$\\
roots $\lambda_j$ & exponent pairs $(\tau_j,\xi_j)$\\
minimal annihilating polynomial & annihilator ideal and its gcd\\
companion matrix & commuting differentiation operators on a finite-dimensional space\\
Lagrange spectral projectors & projectors isolating the individual exponential modes\\
\bottomrule
\end{tabular}
\end{center}
In one variable the roots are isolated points and polynomial factorization gives the whole picture. In two variables the common zero set of polynomials can contain entire curves, so the geometry and the ideal decomposition become more subtle.
\end{explained}

